\chapter{Strong Multiplicity One}\label{chap:smo}

This is the headline chapter. Strong Multiplicity One [Miy, Thm~4.6.12] (equivalently [DS, Thm~5.8.2])
says that a Hecke eigenform is rigidly determined by its eigenvalues at almost all primes: if a
normalised newform $f$ in $S_k(N,\chi)$ and any cusp-form Hecke eigenfunction $g$ in $S_k(N,\chi)$
share the eigenvalues $\lambda_n$ for all $n$ with $(n,N)=1$ outside a finite set, then $g=c\,f$ for a
single constant $c\in\mathbb{C}$. The result rests squarely on the commutativity of the $\mathrm{GL}_2$
Hecke algebra established in the previous chapters (Theorem~\ref{thm:commring-hecke-algebra}) together
with its multiplication table (Theorem~\ref{thm:T-sum-mul}): commutativity is what makes simultaneous
diagonalisation possible, and the rigidity of the eigenvalue systems is read off from the structure of
the $T(n)$. The proof decomposes $g$ along the Petersson-orthogonal new/old splitting, shows the new
part is a multiple of $f$, and shows the old part vanishes.

\section{Eigenforms, newforms, and the new/old decomposition}

\begin{definition}[Nebentypus space $S_k(N,\chi)$]\label{def:modFormCharSpace}
\lean{HeckeRing.GL2.modFormCharSpace}
\leanok
For a positive integer $N$, a weight $k$, and a Dirichlet character $\chi$ modulo $N$, the
\emph{Nebentypus space} (modular-form version $M_k(N,\chi)$, cusp version $S_k(N,\chi)$) is the
common eigenspace
\[
  M_k(N,\chi) \;=\; \bigcap_{d} \ker\bigl(\langle d\rangle - \chi(d)\bigr),
\]
the intersection over all units $d$ modulo $N$ of the $\chi(d)$-eigenspaces of the diamond operators
$\langle d\rangle$. Equivalently, $f$ lies in $M_k(N,\chi)$ if and only if $\langle d\rangle f=\chi(d)\,f$
for every $d$ coprime to $N$.
\end{definition}

\begin{definition}[Eigenform]\label{def:eigenform}
\lean{HeckeRing.GL2.Eigenform}
\uses{def:modFormCharSpace, thm:commring-hecke-algebra, thm:T-sum-mul}
\leanok
An \emph{eigenform} of level $\Gamma_1(N)$ and weight $k$ is a cusp form $f$ carrying a Nebentypus
character $\chi$, so that $f$ lies in the eigenspace $S_k(N,\chi)$, and which is a simultaneous
eigenfunction of all the Hecke operators $T(n)$ with $(n,N)=1$. Because the Hecke operators commute
(Theorem~\ref{thm:commring-hecke-algebra}, whose product law is the multiplication table
Theorem~\ref{thm:T-sum-mul}), such common eigenfunctions exist in abundance: a finite-dimensional space
of cusp forms is simultaneously diagonalised by the commuting family $\{T(n)\}_{(n,N)=1}$. The eigenvalue
is recorded as $T(n)f=\lambda_n\,f$, where $\lambda_n=\chi(n)\,\nu_n$ relates the classical Hecke
eigenvalue $\lambda_n$ to the underlying Hecke-ring eigenvalue $\nu_n$ through the diamond factor $\chi(n)$.
\end{definition}

\begin{definition}[Old subspace $S_k^{\flat}(N)$]\label{def:cuspFormsOld}
\lean{HeckeRing.GL2.cuspFormsOld}
\leanok
The \emph{old subspace} $S_k^{\flat}(N)$ is the span of all level-raising images: for each proper divisor
$M\mid N$ (with $M\neq N$) and each divisor $\ell$ of $N/M$ with $\ell>1$, the degeneracy map
$f\mapsto f(\ell z)$ carries the cusp space at level $M$ into level $N$, and $S_k^{\flat}(N)$ is the
subspace generated by all such images,
\[
  S_k^{\flat}(N) \;=\; \operatorname{span}\bigl\{\, f(\ell z) \ :\ f\in S_k(\Gamma_1(M)),\ M\mid N,\ M\neq N \,\bigr\}.
\]
This is the [DS, (5.18)] definition of the forms ``coming from lower level''.
\end{definition}

\begin{definition}[New subspace $S_k^{\sharp}(N)$]\label{def:cuspFormsNew}
\lean{HeckeRing.GL2.cuspFormsNew}
\uses{def:cuspFormsOld}
\leanok
The \emph{new subspace} $S_k^{\sharp}(N)$ is the orthogonal complement of the old subspace
$S_k^{\flat}(N)$ inside $S_k(\Gamma_1(N))$ with respect to the Petersson inner product,
\[
  S_k^{\sharp}(N) \;=\; \bigl(S_k^{\flat}(N)\bigr)^{\perp},
\]
following [DS, (5.19)]. Concretely, a cusp form lies in $S_k^{\sharp}(N)$ exactly when it is Petersson-orthogonal
to every oldform. The Petersson form is positive definite on cusp forms, so $S_k^{\sharp}(N)\cap S_k^{\flat}(N)=0$.
\end{definition}

\begin{definition}[Old part]\label{def:oldPart}
\lean{HeckeRing.GL2.oldPart}
\uses{def:cuspFormsOld, def:cuspFormsNew}
\leanok
Since $S_k(\Gamma_1(N))=S_k^{\flat}(N)\oplus S_k^{\sharp}(N)$ as an orthogonal direct sum, every cusp form
$f$ has a unique decomposition into an old and a new component. The \emph{old part} of $f$ is its image
under the projection of $S_k(\Gamma_1(N))$ onto $S_k^{\flat}(N)$ along $S_k^{\sharp}(N)$.
\end{definition}

\begin{definition}[New part]\label{def:newPart}
\lean{HeckeRing.GL2.newPart}
\uses{def:cuspFormsOld, def:cuspFormsNew}
\leanok
The \emph{new part} of a cusp form $f$ is its image under the projection of $S_k(\Gamma_1(N))$ onto the
new subspace $S_k^{\sharp}(N)$ along the old subspace $S_k^{\flat}(N)$. Together with the old part it
recovers $f$: one has $f=\operatorname{old}(f)+\operatorname{new}(f)$ with $\operatorname{old}(f)\in S_k^{\flat}(N)$
and $\operatorname{new}(f)\in S_k^{\sharp}(N)$.
\end{definition}

\begin{definition}[Newform]\label{def:newform}
\lean{HeckeRing.GL2.Newform}
\uses{def:eigenform, def:cuspFormsNew}
\leanok
A \emph{newform} (normalised newform) of level $\Gamma_1(N)$ and weight $k$ is an eigenform $f$ that lies
in the new subspace $S_k^{\sharp}(N)$ and is normalised so that its first Fourier coefficient is $a_1=1$.
By Atkin--Lehner uniqueness [DS, Thm~5.8.2] a newform is determined by its Hecke eigenvalues away from the
level; Theorem~\ref{thm:smo-clean} below is exactly this rigidity in the present formalisation.
\end{definition}

\section{Coefficients, leading terms, and the eigenform splitting}

\begin{theorem}[Miyake 4.5.15(1): $a_n=a_1\lambda_n$]\label{thm:coeff-eq-a1-lambda}
\lean{HeckeRing.GL2.Eigenform.coeff_eq_coeff_one_mul_eigenvalue}
\uses{def:eigenform, def:modFormCharSpace}
\leanok
Let $g$ be an eigenform lying in the Nebentypus space $S_k(N,\chi)$, and let $n$ be coprime to $N$. Then
the Fourier coefficients of $g$ satisfy
\[
  a_n(g) \;=\; a_1(g)\,\lambda_n(g),
\]
where $\lambda_n(g)$ is the classical Hecke eigenvalue of $g$ at $n$.
\end{theorem}

\begin{proof}
\uses{def:eigenform, def:modFormCharSpace}
\leanok
This is [Miy, Lemma~4.5.15(1)] in its un-normalised form. Apply the Hecke operator $T(n)$ to $g$. On the
one hand the eigenform relation gives $T(n)g=\lambda_n(g)\,g$, so the first Fourier coefficient of $T(n)g$
is $\lambda_n(g)\,a_1(g)$. On the other hand the explicit formula for the Fourier coefficients of $T(n)g$
expresses $a_m(T(n)g)$ as a sum over the common divisors of $m$ and $n$; taking $m=1$ and using $(n,N)=1$
the sum collapses to its single $d=1$ term, which is precisely $a_n(g)$. Equating the two expressions for
the first coefficient of $T(n)g$ yields $a_n(g)=a_1(g)\,\lambda_n(g)$.
\end{proof}

\begin{theorem}[Miyake 4.6.11: $a_1\neq 0$ for new eigenforms]\label{thm:a1-ne-zero}
\lean{HeckeRing.GL2.coeff_one_ne_zero_of_mem_cuspFormsNew_of_eigen_unconditional}
\uses{def:eigenform, def:cuspFormsNew, thm:coeff-eq-a1-lambda}
\leanok
Let $g$ be a nonzero eigenform lying in $S_k(N,\chi)\cap S_k^{\sharp}(N)$, that is, a nonzero common
eigenfunction of the $T(n)$, $(n,N)=1$, in the new subspace. Then its leading Fourier coefficient is
nonzero, $a_1(g)\neq 0$.
\end{theorem}

\begin{proof}
\uses{def:eigenform, def:cuspFormsNew, thm:coeff-eq-a1-lambda}
\leanok
This is [Miy, Lemma~4.6.11]. Argue by contraposition: suppose $a_1(g)=0$. By
Theorem~\ref{thm:coeff-eq-a1-lambda} every coprime coefficient then vanishes,
$a_n(g)=a_1(g)\,\lambda_n(g)=0$ for all $(n,N)=1$, and the zeroth coefficient vanishes because $g$ is
cuspidal. A form whose Fourier coefficients all vanish away from the level necessarily lies in the old
subspace (this is the content of [Miy, Thm~4.6.8] in the per-character formulation), so $g\in S_k^{\flat}(N)$.
But $g$ also lies in the new subspace $S_k^{\sharp}(N)$, and the Petersson form is positive definite, so
$S_k^{\flat}(N)\cap S_k^{\sharp}(N)=0$ forces $g=0$, contradicting the hypothesis. Hence $a_1(g)\neq 0$.
\end{proof}

\begin{theorem}[Eigenform old/new splitting]\label{thm:eigenform-decomp}
\lean{HeckeRing.GL2.exists_eigenform_decomposition_mem_cuspFormsNew}
\uses{def:eigenform, def:cuspFormsNew, def:oldPart, def:newPart}
\leanok
Let $g$ be a cusp form in the new subspace $S_k^{\sharp}(N)$ whose underlying modular form lies in the
Nebentypus space $S_k(N,\chi)$. Then $g$ is a finite sum of common Hecke eigenforms,
\[
  g \;=\; \sum_{i} h_i,
\]
each summand $h_i$ again lying in $S_k(N,\chi)\cap S_k^{\sharp}(N)$ and being a common eigenfunction of all
$T(n)$ with $(n,N)=1$.
\end{theorem}

\begin{proof}
\uses{def:eigenform, def:cuspFormsNew, def:oldPart, def:newPart}
\leanok
The new subspace is stable under every Hecke operator $T(n)$ with $(n,N)=1$ (DS Prop~5.6.2), and so is the
Nebentypus eigenspace $S_k(N,\chi)$. Since the operators $\{T(n)\}_{(n,N)=1}$ form a commuting family of
diagonalisable operators on the finite-dimensional space $S_k(N,\chi)$
(Theorem~\ref{thm:commring-hecke-algebra}), the simultaneous spectral decomposition of $S_k(N,\chi)$ into
joint eigenspaces restricts to the invariant subspace $S_k^{\sharp}(N)\cap S_k(N,\chi)$. Writing $g$ in
this joint eigenbasis exhibits it as a finite sum of common eigenforms, each of which is itself new and of
Nebentypus $\chi$.
\end{proof}

\section{The two halves of Theorem 4.6.12}

\begin{theorem}[New part is a multiple of $f$]\label{thm:newPart-eq-smul}
\lean{HeckeRing.GL2.newPart_eq_smul_of_shared_eigenvalues}
\uses{def:newform, def:newPart, thm:coeff-eq-a1-lambda, thm:a1-ne-zero, thm:eigenform-decomp}
\leanok
Let $f$ be a normalised newform in $S_k(N,\chi)$ and let $g^{(1)}$ be a common Hecke eigenfunction in the
new subspace lying in $S_k(N,\chi)$ that shares $f$'s eigenvalues $\lambda_n$ for all $(n,N)=1$ outside a
finite set $S$. Then
\[
  g^{(1)} \;=\; b_1\,f, \qquad b_1=a_1(g^{(1)}).
\]
\end{theorem}

\begin{proof}
\uses{def:newform, def:newPart, thm:coeff-eq-a1-lambda, thm:a1-ne-zero, thm:eigenform-decomp}
\leanok
This is the new-part half of [Miy, Thm~4.6.12]. If $g^{(1)}=0$ then $b_1=0$ and the identity is trivial,
so assume $g^{(1)}\neq 0$; then $b_1\neq 0$ by Theorem~\ref{thm:a1-ne-zero}. Sharing eigenvalues off the
finite set $S$ in fact forces equality of \emph{all} coprime eigenvalues (a nonzero coprime eigenvalue and
the multiplicativity of the eigenvalue system propagate the agreement to every $(n,N)=1$). Consider
$h=b_1^{-1}g^{(1)}-f$. For $(n,N)=1$ its $n$-th coefficient is
$b_1^{-1}a_n(g^{(1)})-a_n(f)=\lambda_n-\lambda_n=0$, using Theorem~\ref{thm:coeff-eq-a1-lambda} together
with the normalisation $a_1(f)=1$ which gives $a_n(f)=\lambda_n$. A form with vanishing coprime
coefficients lies in the old subspace by [Miy, Thm~4.6.8], so $h\in S_k^{\flat}(N)$; but $h$ is a
difference of new forms, hence $h\in S_k^{\sharp}(N)$. Orthogonality of the new and old subspaces forces
$h=0$, i.e.\ $b_1^{-1}g^{(1)}=f$, which is the claim.
\end{proof}

\begin{theorem}[Old part vanishes]\label{thm:oldPart-eq-zero}
\lean{HeckeRing.GL2.oldPart_eq_zero_of_shared_eigenvalues}
\uses{def:newform, def:oldPart, thm:a1-ne-zero, thm:eigenform-decomp}
\leanok
Let $f$ be a nonzero newform in $S_k(N,\chi)$ and let $g^{(2)}$ be a common Hecke eigenfunction in the
($\chi$-refined) old subspace, lying in $S_k(N,\chi)$, that shares $f$'s eigenvalues $\lambda_n$ for all
$(n,N)=1$ outside a finite set $S$. Then $g^{(2)}=0$.
\end{theorem}

\begin{proof}
\uses{def:newform, def:oldPart, thm:a1-ne-zero, thm:eigenform-decomp}
\leanok
This is steps (i)--(ii) of [Miy, Thm~4.6.12]. Suppose $g^{(2)}\neq 0$. By the description of the old
subspace [Miy, Lemma~4.6.9] one may write $g^{(2)}$ as a sum of level-raises $h_v(\ell_v z)$ of newforms
$h_v$ at proper divisors $M_v$ of $N$; by Theorem~\ref{thm:eigenform-decomp} the $h_v$ may be taken to be
common eigenforms, and by [Miy, Lemma~4.6.2] their level-raises remain eigenforms with the same eigenvalues.
Eigenfunctions with distinct eigenvalues are Petersson-orthogonal, hence linearly independent, so the
components whose eigenvalues differ from those of $f$ contribute nothing; removing them leaves a nonzero
new eigenform $h$ at a proper divisor $M$ whose eigenvalues match $f$'s off $S$. By
Theorem~\ref{thm:a1-ne-zero} its leading coefficient $c_1'$ is nonzero. The difference $h-c_1'f$ has
vanishing coprime coefficients (Theorem~\ref{thm:coeff-eq-a1-lambda} and $a_1(f)=1$), so it lies in the old
subspace by [Miy, Thm~4.6.8]; and $h$ itself, being a form essentially of level $M\neq N$, also lies in the
old subspace. Therefore
\[
  f \;=\; c_1'^{-1}h - c_1'^{-1}(h-c_1'f)
\]
exhibits $f$ as an element of the old subspace. This contradicts $f$ being a nonzero newform (the new and
old subspaces meet only in $0$). Hence $g^{(2)}=0$.
\end{proof}

\section{Strong Multiplicity One}

\begin{theorem}[Strong Multiplicity One, Miyake 4.6.12]\label{thm:smo-constMul}
\lean{HeckeRing.GL2.strongMultiplicityOne_constMul}
\uses{def:newform, def:eigenform, def:oldPart, def:newPart, thm:newPart-eq-smul, thm:oldPart-eq-zero}
\leanok
Let $f$ be a normalised newform in $S_k(N,\chi)$, and let $g$ be any cusp-form Hecke eigenfunction in
$S_k(N,\chi)$. If $f$ and $g$ are common eigenfunctions of $T(n)$ with the same eigenvalue $\lambda_n$ for
every $n$ with $(n,N)=1$ outside a finite set $S$, then $g$ is a constant multiple of $f$:
\[
  g \;=\; c\,f \qquad\text{for some } c\in\mathbb{C}.
\]
\end{theorem}

\begin{proof}
\uses{def:newform, def:eigenform, def:oldPart, def:newPart, thm:newPart-eq-smul, thm:oldPart-eq-zero}
\leanok
This is [Miy, Thm~4.6.12]. Normalise $a_1(f)=1$, which is legitimate since $a_1(f)\neq 0$ by
[Miy, Lemma~4.6.11], and split $g$ along the orthogonal new/old decomposition,
$g=g^{(1)}+g^{(2)}$ with $g^{(1)}\in S_k^{\sharp}(N)$ and $g^{(2)}\in S_k^{\flat}(N)$. Because both the new
subspace and the old subspace are stable under the Hecke operators (DS Prop~5.6.2 / [Miy, Lemma~4.6.10]),
each of $g^{(1)}$ and $g^{(2)}$ is again a common eigenfunction of $T(n)$ sharing $f$'s eigenvalue
$\lambda_n$ for $(n,N)=1$ outside $S$.

For the new part, Theorem~\ref{thm:newPart-eq-smul} gives $g^{(1)}=b_1\,f$ where $b_1=a_1(g^{(1)})$: the
difference $b_1^{-1}g^{(1)}-f$ has vanishing coprime coefficients by [Miy, Lemma~4.5.15(1)] and therefore
lies in the old subspace by [Miy, Thm~4.6.8], while being new, hence is $0$. For the old part,
Theorem~\ref{thm:oldPart-eq-zero} gives $g^{(2)}=0$: assuming otherwise, one descends $g^{(2)}$ to a
nonzero new eigenform $h$ at a proper divisor sharing $f$'s eigenvalues ([Miy, Lemma~4.6.9],
[Miy, Lemma~4.6.2], and linear independence of distinct-eigenvalue eigenforms), and then
$h-c_1'f\in S_k^{\flat}(N)$ together with $h\in S_k^{\flat}(N)$ forces $f\in S_k^{\flat}(N)$, contradicting
that $f$ is a nonzero newform. Combining the two halves, $g=g^{(1)}+g^{(2)}=b_1\,f$, so $g=c\,f$ with
$c=b_1=a_1(g)$.
\end{proof}

\begin{theorem}[Multiplicity One for newforms, DS 5.8.2]\label{thm:smo-clean}
\lean{HeckeRing.GL2.strongMultiplicityOne_axiom_clean_unconditional}
\uses{thm:smo-constMul, def:newform}
\leanok
Let $f$ and $g$ be normalised newforms in the same Nebentypus space $S_k(N,\chi)$. If $f$ and $g$ have the
same Hecke eigenvalue $\lambda_n$ for every $n$ with $(n,N)=1$ outside a finite set $S$, then $f=g$.
\end{theorem}

\begin{proof}
\uses{thm:smo-constMul, def:newform}
\leanok
This is [DS, Thm~5.8.2] and follows immediately from the headline Theorem~\ref{thm:smo-constMul}. Applying
it to the newform $f$ and the eigenfunction $g$ produces a constant $c$ with $g=c\,f$. Comparing first
Fourier coefficients and using that both forms are normalised, $1=a_1(g)=c\,a_1(f)=c$, so $c=1$ and hence
$g=f$.
\end{proof}
